% Discussion chapter: the paper's discussion and limitations are input, then thesis-only material.
\section{Discussion}
\label{sec:discussion}
% Written from the final tables (results/final/stats_held.json, results/bdcxr/stats.json). Failed hypotheses are named first.
\paragraph{Principal findings.} Three pre-specified results did not hold: the soft-tissue image alone did not improve any downstream task
(Table~\ref{tab:downstream_cls}, Table~\ref{tab:downstream_det}), SSIM fell short of the reference model by more than its margin
(Table~\ref{tab:h2}), and conformal intervals calibrated on synthetic data covered 14\,\% instead of 90\,\% of real bone pixels
(Table~\ref{tab:reliability}); recalibrated on 24 real pairs in an exploratory analysis, they covered 89\,\% (Table~\ref{tab:bone_recal}). Real-nodule contrast was mixed: higher than with the reference model on JSRT, lower on Node21
(Table~\ref{tab:nodule_cnr}). The central results held. Lesion absorption fell steadily with the twin weight (Table~\ref{tab:dose_response}); the
twin ensemble absorbed a third of what the released CT-supervised model absorbed and less than a third of what the matched control absorbed, on
an audit set no model had seen during development (Table~\ref{tab:audit_contrasts}); and it did so at non-inferior bone-region fidelity. The effect
replicated with two other encoders (Table~\ref{tab:encoders}) and on paediatric radiographs from a population no model was trained on (Section~\ref{sec:bdcxr}).

\paragraph{What the audit adds.} Fidelity metrics and lesion absorption rank suppressors differently (Figure~\ref{fig:pareto}). The v1 model has
the highest bone-region PSNR of the CT-supervised models and the highest absorption; ours and the reference model are within 0.06\,dB of each
other and a factor of three apart in absorption. BS-LDM, trained on dual-energy pairs from the same hospital, is about 7\,dB ahead in fidelity on
that hospital's images, yet in the development audit it absorbed 0.42 of an inserted lesion (at site positions not paired with the other
methods, so only descriptive). A suppressor that matches the dual-energy target closely can still remove lesions: fidelity to that target does
not protect what the target never showed it. Absorption must nevertheless be read beside suppression strength. The reference model removes more
rib-edge energy than ours (48\,\% against 30\,\%) and more non-rib energy (22\,\% against 8\,\%), so part of its higher absorption may reflect
stronger suppression in general; ours absorbs a third as much while removing about 60\,\% of the rib-edge energy.

\paragraph{What the regulariser costs.} On real paired radiographs the twin costs no measurable bone-region fidelity up to $w=16$ (21.22 to
21.28\,dB across $w=0\ldots16$); at $w=32$ and 64 the PSNR drops by about 0.3\,dB, so $w=16$ is the usable
operating point, chosen on development data before the freeze. The cost appears where the truth is exact: on the Qure-projection test split the
bone error of the twin ensemble is 41\,\% higher than that of the control (MAE 0.0041 against 0.0029), and in the ablation the validation error rises
from 0.0033 to 0.0057. The regulariser therefore trades some synthetic bone accuracy for lesion invariance; on real images that trade was not
visible. The ablation shows what carries the effect: gradients through both twin branches and a consistency term confined to the pasted lesion.
Showing lesions during training without the twin term (CLIA) did not lower absorption.

\paragraph{Relation to other methods.} The reference model is a SwinV2-B at $1024$~px trained with a far larger sample budget; we compare its
released weights as a user would run them, and we cannot say how much of the absorption difference is due to our regulariser rather than to
training differences, except through the matched no-twin control, which absorbs slightly more than the reference model (0.178 against 0.149). Diffusion models trained on real
dual-energy pairs \citep{sun2026bsldm,sun2026debonedit} answer a different question: they need paired data from the target site, sample slowly
(about 2.5 minutes per image for BS-LDM on a T4 GPU) and were not designed for lesion invariance.

\paragraph{Advice for evaluation and use.} We suggest that a bone suppressor be audited for lesion absorption before it is used, with the audit
reported next to a suppression measure; that suppressed images be offered beside the original rather than instead of it, which is also the only
use in which they helped downstream models here; and that the ensemble spread be used as a warning flag rather than as a calibrated interval: it
rose 2.7-fold on paediatric radiographs, but intervals calibrated on synthetic data did not cover real errors, and recalibration needed real
pairs. Whether classifiers that combine the radiograph with the soft-tissue and bone images in separate streams gain from suppression is tested
in \clsstudy.

\paragraph{Future work.} Prediction at $1024$~px, combining CT-derived and real dual-energy supervision with the twin term, conformal
recalibration on real pairs from more sites, lesion types other than nodules (consolidation, as in the paediatric pneumonia set), paediatric training data, and a
reader study that measures detectability rather than signal.


\section{Limitations}
\label{sec:limitations}
% PLAN (about 0.5 page, one short paragraph per item). Items 1-9 are known now from DEVIATIONS.md, DECISIONS.md and SAP_v2.md; add what the final
% results reveal. State each limitation with its consequence for how the results must be read. Do not soften a limitation that changes a conclusion.
\begin{enumerate}
\item \emph{Synthetic supervision and domain gap.} Training uses rendered or translated projections; the bone image of a real radiograph is never
 observed. Real-image fidelity is measured against dual-energy soft-tissue images, which differ from the radiograph in tone and in high-frequency
 content, so raw and band-limited metrics are both reported and the calibrated metrics are not absolute.
\item \emph{One source of paired real data.} SZCH-X-Rays comes from one partner hospital and one acquisition set-up (dual-exposure DES on a GE Discovery XR656).
\item \emph{The audit is linear and additive.} Lesions are added in image units to a real radiograph; beam hardening, scatter and the tone mapping of
 the real image are not modelled, so $A$ measures the leak of an added signal and not clinical detectability. $A=0$ is reached by doing nothing, and
 $A$ is therefore read next to a suppression gain.
\item \emph{Shared insertion mechanism.} The regulariser and the audit both insert lesions. The confirmatory audit, the real-nodule contrast and the
 downstream subgroup analysis were added to limit this, but they do not remove it.
\item \emph{Development use of SZCH.} All 75 SZCH images were used in development (audit) and 24 for paired decisions; only the 51 held-out pairs are
 clean for fidelity claims, and the SZCH audit is development evidence only. The plan changed once after the SZCH pairs were looked at
 (SAP\_v2.md, DEVIATIONS.md).
\item \emph{Resolution and backbone.} Our models predict at 512 px with tiny encoders (ConvNeXt-V2-T for ours, SwinV2-T for the dose-response);
 the reference model predicts at 1024 px with SwinV2-B. The SSIM gap sits in fine texture: with a band limit of $\sigma=4$ it shrinks from 0.021 to
 0.002, and ours keeps about 9\,\% more high-frequency energy than the dual-energy target in non-bone regions (the reference model 1\,\% less). Whether
 prediction at 1024 px closes the remaining gap was not tested.
\item \emph{Operating point.} The twin weight trades absorption against fidelity and against bone accuracy on exact-truth data (Section~\ref{sec:discussion}).
 Only $w=16$ was evaluated with the full protocol; the dose-response shows that $w=32$ and 64 lower bone-region PSNR by about 0.3\,dB.
\item \emph{Uncertainty.} Conformal calibration treats pixels as exchangeable, so coverage is marginal over pixels, and it is calibrated on synthetic
 data; coverage on real images is reported as measured under shift. The recalibration on real pairs was chosen after the shift was seen and
 used 24 pairs from one hospital. Ensemble members share training data.
\item \emph{Scope.} Training and paired evaluation use adult frontal radiographs; paediatric images enter only the external validation, where no
 paired truth exists; no lateral images and no reader study. Research use only and not a medical device. The weights are non-commercial because
 the training data are.
\item \emph{Overlap maps.} On JSRT the bone overlap of each audit site came from the reference model's bone map, not from the consensus map the
 plan names, so the JSRT interaction is defined by a comparator; on VinDr-RibCXR it comes from radiologist rib masks and shows no interaction.
\item \emph{Downstream and external validation.} Downstream models were trained once per seed with fixed settings, not tuned per input; the
 subgroup analysis by bone overlap is not estimable on TBX11K, where every test box overlaps bone. BDCXR-3257 is a single-centre set with
 image-level labels and no patient identifiers, so repeated images of one child can only be caught by the near-duplicate grouping; its audit
 inserts adult nodules into children's radiographs.
\item \emph{Post-freeze corrections.} Several analysis errors found in an independent check after the first locked run were corrected
 (a resampling error that narrowed interaction intervals, incomplete downstream statistics, and detection metrics that depended on the order of
 tied scores by up to 0.00055); some evaluations ran on cloud GPUs whose floating-point results differ from the local ones by at most
 0.0024\,dB, with the audit sites pinned to the local lung masks. No pre-specified verdict changed; every change is listed in the deviations log.
\item \emph{Registration.} The analysis plans were hash-frozen in the project's version control before testing but not deposited in a public
 registry.
\end{enumerate}


\section{Threats to validity}
\paragraph{Internal validity.} The operating point, the training source and the encoder were chosen on development data: the 24 SZCH
development pairs, the audit on all 75 SZCH images and synthetic validation data. The confirmatory sets (VinDr-RibCXR, JSRT, the 51 held-out
SZCH pairs, the downstream test sets and BDCXR-3257) were evaluated only after the plans were frozen, and the screening tools that had printed
held-out SZCH numbers next to development numbers were changed and the episode disclosed. Seeds were fixed per arm, and every final contrast
pairs models on identical images and sites; where cloud and local GPUs were mixed, the audit sites were pinned to the same lung masks and the
remaining floating-point difference was measured (at most 0.0024\,dB). The selection of ConvNeXt-V2-T over RAD-DINO was made by a pre-declared
tie-break at two seeds and was not reopened when the third seeds moved the gap just outside the band.

\paragraph{External validity.} Paired real truth comes from one hospital and one dual-energy set-up. Training data are adult CT projections,
and the external validation covers one paediatric hospital. The audit uses nodules only; consolidations, masses, effusions and lines were not
inserted. Lateral, portable and supine radiographs were not studied.

\paragraph{Construct validity.} Absorption measures how much of an added, linear lesion signal moves into the bone estimate. It is not
clinical detectability, and a model that suppresses less absorbs less for that reason alone; it is therefore read with the rib benchmark and
bone-region fidelity. Real-nodule contrast and downstream performance were added because they do not use insertion, and they agree with the
audit against the matched control but not uniformly against the reference model.

\paragraph{Statistical conclusion validity.} Holm correction is applied to the $p$-values within each hypothesis family; confidence intervals
are unadjusted. Bootstrap intervals resample whole images, patients or near-duplicate groups as the plans require; per-class DeLong tests on
NIH ignore the clustering of images within patients, which the macro-level patient bootstrap does not. An independent check after the first
locked run found and fixed several analysis errors, among them a resampling error that had narrowed the interaction intervals and detection
metrics that depended on the order of tied scores; no pre-specified verdict changed, and every correction is listed in
Appendix~\ref{app:deviations}.

\section{Future work}
Five directions follow from the results. Prediction at $1024$~px would show whether the remaining SSIM gap to the reference model is a matter of
resolution. Combining CT-derived supervision with real dual-energy pairs, where available, could raise fidelity without giving up the twin
term, and real pairs from the target site would allow conformal calibration that holds on real images. The audit should be extended to
lesion types other than nodules, such as consolidation, which the paediatric pneumonia set points to. Paediatric training data would test
whether the higher absorption and uncertainty on children can be reduced. Finally, a reader study would measure what the audit cannot:
whether radiologists find more lesions on images from a suppressor with low absorption.
